\ExerciseSection

\begin{enumerate}
\def\labelenumi{\arabic{enumi})}
\item
  Let X be a Lindelöf space (Chapter I, § 9, Exercise 15), let H be a directed set (with respect to the relation \(\geqslant\) ) of continuous real-valued functions on X such that the lower envelope g of the functions of H is continuous. Show that there exists a decreasing sequence of functions \((f_{n})\) belonging to H which converges pointwise to g.
\item
  Let I be a compact interval in R and let \((f_{n})\) be a sequence of monotone real-valued functions defined on I which converges pointwise in I to a continuous function g. Show that g is monotone and that the sequence \((f_{n})\) converges uniformly to g in I.
\end{enumerate}

¶ 3) Let \(\Gamma\) be a simply transitive group of homeomorphisms of \(\mathbf{R}\) endowed with the topology of pointwise convergence. For each \(x \in \mathbf{R}\) let \(s_x\) denote the element of \(\Gamma\) such that \(s_x(0) = x\) . Show that the mapping \(x \to s_x\) is a homeomorphism of \(\mathbf{R}\) onto \(\Gamma\) {[}use the fact that if \(s \in \Gamma\) is such that \(x < s(x)\) for some \(x \in \mathbf{R}\) , then \(y < s(y)\) for all \(y \in \mathbf{R}]\) ; hence show that \(\Gamma\) is isomorphic to \(\mathbf{R}\) as a topological group {[}use Exercise 17 a) of § 3 and Theorem 1 of Chapter V, § 3{]}.

\begin{enumerate}
\def\labelenumi{\arabic{enumi})}
\setcounter{enumi}{3}
\item
  Let X be a compact space and let H be a vector subspace of C(X; R) such that \(|u| \in H\) whenever \(u \in H\) . Then a continuous real-valued function f on X can be uniformly approximated by functions belonging to H if and only if, for each pair x, y of points of X, the function f satisfies every linear relation \(\alpha f(x) = \beta f(y)\) where \((\alpha, \beta)\) are such that \(\alpha\beta \geqslant 0\) and \(\alpha g(x) = \beta g(y)\) for every function \(g \in H\) . {[}Apply Proposition 2 of no. 1, observing that the image of H under the mapping \(u \to (u(x) u(y))\) is the whole plane \(R^{2}\) , or a line of equation \(\alpha X = \beta Y\) with \(\alpha\beta \geqslant 0\) , or else consists of the single point \((0, 0)\) .{]}
\item
  Let X be a compact space and let H be a set of continuous real-valued functions on X. Let R be the equivalence relation `` \(u(x) = u(y)\) for all \(u \in H\)'' . A continuous real-valued function f on X can then be uniformly approximated by polynomials (resp. polynomials with no constant term) in the functions of H if and only if f is constant on each equivalence class mod R {[}resp. if and only if f is constant on each equivalence class mod R and vanishes at all points x such that \(u(x) = 0\) for all \(u \in H\) {]}.
\end{enumerate}

¶ 6) Let X be a completely regular space and let \(\mathcal{C}^{\infty}(\mathbf{X};\mathbf{R})\) denote the normed subalgebra of \(\mathcal{B}(\mathbf{X};\mathbf{R})\) consisting of all bounded continuous real-valued functions on X. In order that every subalgebra of \(\mathcal{C}^{\infty}(\mathbf{X};\mathbf{R})\) , which separates the points of X and contains the constant functions, should be dense in \(\mathcal{C}^{\infty}(\mathbf{X};\mathbf{R})\) , it is necessary and sufficient that X should be compact. {[}Let \(\beta\mathbf{X}\) be the Stone-Čech compactification of X (Chapter IX, § 1, Exercise 7). Show that if \(\beta\mathbf{X}-\mathbf{X}\) is not empty then there are non-dense subalgebras of \(\mathcal{C}^{\infty}(\mathbf{X};\mathbf{R})\) which separate the points of X and contain the constant functions, by observing that \(\mathcal{C}^{\infty}(\mathbf{X};\mathbf{R})\) may be identified with \(\mathcal{C}(\beta\mathbf{X};\mathbf{R}).]\)

\begin{enumerate}
\def\labelenumi{\arabic{enumi})}
\setcounter{enumi}{6}
\tightlist
\item
  Let X be a non-compact, completely regular space.
\end{enumerate}

\begin{enumerate}
\def\labelenumi{\alph{enumi})}
\tightlist
\item
  Show that the following properties are equivalent:
\end{enumerate}

\(\alpha)\) The subalgebra of \(\mathcal{C}^{\infty}(\mathbf{X};\mathbf{R})\) consisting of functions of the form \(c + f\) , where \(c\) is a constant and \(f\) has compact support (Chapter IX, § 4, no. 3), is dense in \(\mathcal{C}^{\infty}(\mathbf{X};\mathbf{R})\) .

\(\beta)\) If \(\beta X\) is the Stone-Čech compactification of \(X\) , then \(\beta X - X\) consists of a single point.

γ) There is only one uniformity on X, compatible with its topology, with respect to which X is precompact.

δ) If A, B are two completely separated closed subsets of X (Chapter IX, § 1, Exercise 11) then one of A, B is compact.

\(\zeta)\) If \(\mathcal{U}\) is any uniformity on \(X\) compatible with its topology and if \(V\) is any entourage of \(\mathcal{U}\) , there is a subset \(L\) of \(X\) whose complement is compact, such that \(L \times L \subset V\) .

θ) There is only one uniformity on X compatible with its topology. {[}Note that α) implies that X is locally compact and therefore open in βX, and hence show that α) → β); to prove that β) → α) use the Weierstrass-Stone theorem. To show that β) → γ), use the fact that the uniformity induced on X by that of βX is the finest uniformity compatible with the topology of X, with respect to which X is precompact. To show that γ) → δ), use Chapter IX, § 1, Exercise 11. To show that δ) → ζ), show first that if V is any entourage of U, there exists z ∈ X such that V(z) is not compact; otherwise X would be complete with respect to U and paracompact (Chapter II, § 4, Exercise 9), therefore normal; show then with the help of δ) that X would be countably compact (Chapter IX, § 2, Exercise 14), and finally use Chapter

II, § 4, Exercise 6. Then apply this result to an entourage V defined by \(f(x, x') \leqslant 1\) , where \(f\) is a continuous pseudometric on X, and deduce from \(\delta\) ) that there is a subset L of X, whose complement is compact, such that \(L \times L \subset V\) . Finally, to establish that \(\zeta) \Longrightarrow \theta\) ), observe that every ultrafilter on X either converges or is finer than the filter of complements of relatively compact subsets of X.{]}

\begin{enumerate}
\def\labelenumi{\alph{enumi})}
\setcounter{enumi}{1}
\item
  Show that a space X which satisfies the equivalent conditions of a) is pseudo-compact (Chapter IX, § 1, Exercise 22). {[}Use Exercise 12 of Chapter IX, § 1, or else observe that if X is not pseudo-compact, there is a sequence \((x_{n})\) of points of X which has no cluster point and a continuous mapping \(f: \mathbf{X} \to [\mathrm{o}, \mathrm{i}]\) such that \(f(x_{2n}) = \mathrm{o}\) and \(f(x_{2n+1}) = \mathrm{i}\) .{]}
\item
  Let Y be a non-compact, completely regular space. Show that the subspace X of \(\beta Y\) which is the complement of a point of \(\beta Y - Y\) satisfies the conditions of a).
\end{enumerate}

\begin{enumerate}
\def\labelenumi{\arabic{enumi})}
\setcounter{enumi}{7}
\tightlist
\item
  Let \((\mathbf{X}_i)_{1\leqslant i\leqslant n}\) be a finite family of compact spaces, let \(\mathbf{X} = \prod_{i=1}^{n}\mathbf{X}_i\) be their product, and for each index \(i\) let \(\mathbf{F}_i\) be a closed subset of \(\mathbf{X}_i\) . Show that every continuous real-valued function on \(\mathbf{X}\) which vanishes on the union of the sets \(\mathbf{F}_i\times \prod_{j\neq i}\mathbf{X}_j\) ( \(1\leqslant i\leqslant n\) ) can be uniformly approximated by finite sums of functions of the form
\end{enumerate}

\[
(x _ {1}, \dots , x _ {n}) \rightarrow u _ {1} (x _ {1}) \dots u _ {n} (x _ {n}),
\]

where \(u_{i}\) is a continuous real-valued function on \(\mathbf{X}_i\) which vanishes on \(\mathbf{F}_i\) ( \(1 \leqslant i \leqslant n\) ).

* 9) Let \((r_n)_{n \geqslant 1}\) be the sequence of distinct rational numbers belonging to the interval \(\mathbf{I} = [\mathbf{o}, \mathbf{i}]\) of \(\mathbf{R}\) , arranged in some order. Define by induction a sequence of closed intervals \(\mathbf{I}_n \subset \mathbf{I}\) , as follows: \(\mathbf{I}_n\) has as its midpoint the point \(r_{k_n}\) with the smallest index not contained in \(\bigcup_{p < n} \mathbf{I}_p\) ; its length is \(\leqslant \mathbf{i}/4^n\) and it meets no interval \(\mathbf{I}_p\) for which \(p < n\) ; \((\mathbf{I}_n)\) is thus a sequence of mutually disjoint closed intervals. On the product space \(\mathbf{I} \times \mathbf{R}\) we define a continuous real-valued function \(u\) as follows: for each integer \(n \geqslant \mathbf{i}\) , the function \(x \to u(x, n)\) is equal to \(\mathbf{i}\) at an interior point of \(\mathbf{I}_n\) , is equal to \(\mathbf{o}\) on the exterior of \(\mathbf{I}_n\) and takes its values in \([\mathbf{o}, \mathbf{i}]\) ; on the other hand, for each \(x \in \mathbf{I}\) , the function \(y \to u(x, y)\) is affine-linear in each of the intervals \([n, n + \mathbf{i}]\) . Show that \(u\) cannot be uniformly approximated in \(\mathbf{I} \times \mathbf{R}\) by linear combinations of functions of the form \(v(x)w(y)\) , where \(v\) is continuous on \(\mathbf{I}\) and \(w\) is continuous and bounded on \(\mathbf{R}\) . {[}In the space \(\mathcal{C}^\infty(\mathbf{R}; \mathbf{R})\) of continuous bounded real-valued functions on \(\mathbf{R}\) , consider the set of partial functions \(y \to u(x, y)\) as \(x\) runs through \(\mathbf{I}\) ; show that there exists an infinite sequence \((u_n)\) of these functions such that \(||u_n|| = \mathbf{I}\) and \(||u_m - u_n|| = \mathbf{I}\) whenever \(m \neq n\) . Deduce that there cannot exist any finite-dimensional vector subspace \(\mathbf{E}\) of \(\mathcal{C}^\infty(\mathbf{R}; \mathbf{R})\) such that the distance of each \(u_n\) from \(\mathbf{E}\) is \(\leqslant 1/4\) ; for otherwise there would exist a sequence \((v_n)\) of points of \(\mathbf{E}\) such that \(||v_n|| = 2\) and \(||v_n - v_m|| \geqslant 1/2\) whenever \(m \neq n\) , contrary to the fact that every finite-dimensional subspace of \(\mathcal{C}^\infty(\mathbf{R}; \mathbf{R})\) is locally compact.{]} *

\begin{enumerate}
\def\labelenumi{\arabic{enumi})}
\setcounter{enumi}{9}
\tightlist
\item
  Use the Weierstrass-Stone theorem to give a new proof of Urysohn's theorem (Chapter IX, § 4, no. 2, Theorem 2) for closed subspaces of a compact space X. {[}If F is a closed subspace of X, consider the set H of restrictions to F of continuous mappings \(\mathbf{X} \to \mathbf{R}\) , and observe that if \(f \in \mathbf{H}\) and \(|f(x)| \leqslant a\) for all \(x \in \mathbf{F}\) , then there exists a function \(g \in \mathbf{H}\) which coincides with \(f\) on F and is such that \(|g(x)| \leqslant a\) in X.{]}
\end{enumerate}

¶ 11) a) Let X be a completely regular space and let Y be a closed subspace of X; the mapping \(\varphi: u \to u|\mathbf{Y}\) of \(\mathcal{C}^{\infty}(\mathbf{X}; \mathbf{R})\) into \(\mathcal{C}^{\infty}(\mathbf{Y}; \mathbf{R})\) is a continuous linear mapping of norm \(\leqslant 1\) . If X is normal, \(\varphi\) is a surjective strict morphism; and if \(\mathbf{X}_{\mathbf{Y}}\) denotes the quotient space of X obtained by identifying all the points of Y ( \(\mathbf{X}_{\mathbf{Y}}\) is normal by Exercise 15 of Chapter IX, § 4), then the kernel of \(\varphi\) can be identified (as a normed space) with the subspace of \(\mathcal{C}^{\infty}(\mathbf{X}_{\mathbf{Y}}; \mathbf{R})\) consisting of functions which vanish at \(\omega\) , the canonical image of Y in \(\mathbf{X}_{\mathbf{Y}}\) .

\begin{enumerate}
\def\labelenumi{\alph{enumi})}
\setcounter{enumi}{1}
\tightlist
\item
  Suppose that X is metrizable and that the frontier of Y in X is of countable type. Show that \(\varphi: u \to u|Y\) is a strict morphism which has a right inverse (Chapter III, § 6, no. 2, Proposition 3). {[}Let \(d\) be a metric compatible with the topology of X; let \((a_n)\) be a sequence of points of Fr (Y), dense in Fr (Y), and let \(\mathbf{V}_{n,m}\) be the set of points \(x \in \mathbf{X}\) such that \(d(x, a_n) < \mathrm{i}/m\) and \(d(x, \mathbf{Y}) > \mathrm{i}/2m\) ; construct a family of continuous mappings \(f_{n,m}: \mathbf{X} \to [\mathrm{o}, \mathrm{i}]\) such that the support of \(f_{n,m}\) is contained in \(\mathbf{V}_{n,m}\) and such that \(\sum_{n,m} f_{n,m}(x) = \mathrm{i}\) for all \(x \in \complement Y\) , the family \((f_{n,m})\) being uniformly summable in some neighbourhood of each point of \(\complement Y\) . Show then that, for every \(v \in \mathcal{C}^\infty(\mathbf{Y}; \mathbf{R})\) , the function \(u\) which coincides with \(v\) on Y and is equal to \(\sum_{n,m} v(a_n) f_{n,m}(x)\) for all \(x \in \complement Y\) belongs to \(\mathcal{C}^\infty(\mathbf{X}; \mathbf{R})\) and is such that \(\varphi(u) = v.\) {]} (*)
\end{enumerate}

\begin{enumerate}
\def\labelenumi{\arabic{enumi})}
\setcounter{enumi}{11}
\tightlist
\item
  \begin{enumerate}
  \def\labelenumii{\alph{enumii})}
  \tightlist
  \item
    Let X be a compact space, and let \(f: \mathbf{X} \to \mathbf{Y}\) be a continuous surjection of X onto a compact space Y. Show that the mapping \(^{a}f: u \to u \circ f\) of \(\mathcal{C}(\mathbf{Y}; \mathbf{R})\) into \(\mathcal{C}(\mathbf{X}; \mathbf{R})\) is a (norm-preserving) isomorphism of \(\mathcal{C}(\mathbf{Y}; \mathbf{R})\) onto a closed subalgebra of \(\mathcal{C}(\mathbf{X}; \mathbf{R})\) containing the identity element.
  \end{enumerate}
\end{enumerate}

\begin{enumerate}
\def\labelenumi{\alph{enumi})}
\setcounter{enumi}{1}
\item
  Conversely, let A be a closed subalgebra of \(\mathcal{C}(\mathbf{X};\mathbf{R})\) containing the identity element. Show that there exists a continuous surjection \(f\) of \(\mathbf{X}\) onto a compact space \(\mathbf{Y}\) such that \(^a f\) is an isomorphism of \(\mathcal{C}(\mathbf{Y};\mathbf{R})\) onto A. {[}If \((u_{\lambda})_{\lambda \in \mathbf{L}}\) is a family which is dense in A, consider the continuous mapping \(x\to (u_{\lambda}(x))\) of \(\mathbf{X}\) into \(\mathbf{R}^{\mathrm{L}}\) and use the Weierstrass-Stone theorem.{]}
\item
  Deduce from b) that, for each sequence \((u_{n})\) of elements of \(\mathcal{C}(\mathbf{X};\mathbf{R})\) , there exists a metrizable compact space Y and a continuous surjective mapping \(f: X \to Y\) such that the \(u_{n}\) belong to the image of \(\mathcal{C}(\mathbf{Y};\mathbf{R})\) under \(^{a}f\) .
\end{enumerate}

¶ 13) Let X be a compact space and let A denote the normed algebra C(X; R). For each subset M of A, let V(M) denote the set of all \(x \in X\) such that \(u(x) = o\) for all \(u \in M\) . For each subset Y of X, let \(\Im(Y)\) denote the set of all \(u \in A\) such that \(u(x) = o\) for all \(x \in Y\) .

\begin{enumerate}
\def\labelenumi{\alph{enumi})}
\tightlist
\item
  V(M) is a closed subspace of X and \(\Im(Y)\) is a closed ideal of A. If a is the ideal of A generated by a subset M of A, then \(\mathrm{V}(\mathrm{M})=\mathrm{V}(\mathfrak{a})\) ; V and \(\Im\) are order-reversing mappings (with respect to the order relation of inclusion), and we have \(V\{o\}=X\) ; \(\mathrm{V}(\mathrm{A})=\varnothing\) ; \(\mathrm{V}(\{u\})=\varnothing\) for every unit u of A;
\end{enumerate}

\[
\mathrm{V} \left(\bigcup_ {\lambda \in \mathrm{L}} \mathrm{M} _ {\lambda}\right) = \mathrm{V} \left(\sum_ {\lambda \in \mathrm{L}} \mathrm{M} _ {\lambda}\right) = \bigcap_ {\lambda \in \mathrm{L}} \mathrm{V} (\mathrm{M} _ {\lambda})
\]

for every family \((\mathbf{M}_{\lambda})_{\lambda \in \mathbf{L}}\) of subsets of A; \(\mathbf{V}(\mathbf{M}. \mathbf{M}') = \mathbf{V}(\mathbf{M}) \cup \mathbf{V}(\mathbf{M}')\) ; \(\Im(\mathbf{X}) = \{\mathrm{o}\}\) ; \(\Im(\emptyset) = \mathbf{A}\) ; \(\Im\left(\bigcup_{\lambda \in \mathbf{L}} \mathbf{Y}_{\lambda}\right) = \bigcap_{\lambda \in \mathbf{L}} \Im(\mathbf{Y}_{\lambda})\) for every family \((\mathbf{Y}_{\lambda})_{\lambda \in \mathbf{L}}\) of subsets of X.

\begin{enumerate}
\def\labelenumi{\alph{enumi})}
\setcounter{enumi}{1}
\item
  Show that if a is any ideal of A, then \(\Im(\mathrm{V}(a)) = \overline{a}\) , and that if Y is any subset of X, then \(\mathrm{V}(\Im(\mathrm{Y})) = \overline{\mathrm{Y}}\) . (For the first of these relations, observe that \(\mathrm{V}(\overline{a}) = \mathrm{V}(a)\) and that we may therefore suppose a closed; note that if x, y are two distinct points of \(\mathrm{X} - \mathrm{V}(a)\) , there exists \(u \in a\) such that \(u(x) \neq u(y)\) , and use Proposition 4 of no. 2. To prove the second relation, use Urysohn's theorem.)
\item
  Deduce from b) that \(x \to \Im(\{x\})\) is a bijection of X onto the set of maximal ideals of A, which are therefore closed. An ideal of A is closed if and only if it is an intersection of maximal ideals. The ring A has no radical.
\item
  Show that the mapping \(e \to \mathrm{V}(\{e\})\) is a bijection of the set of idempotents of A onto the set of subsets of X which are both open and closed.
\end{enumerate}

A point \(x \in \mathbf{X}\) is isolated if and only if the maximal ideal \(\Im(\{x\})\) is principal.

\begin{enumerate}
\def\labelenumi{\arabic{enumi})}
\setcounter{enumi}{13}
\tightlist
\item
  Let \(\mathbf{X}\) be a topological space. For each function \(u \in \mathcal{C}(\mathbf{X}; \mathbf{R})\) such that \(u(x) > 0\) for all \(x \in \mathbf{X}\) , let \(\mathbf{V}_u\) be the set of all \(f \in \mathcal{C}(\mathbf{X}; \mathbf{R})\) such that \(|f| \leqslant u\) .
\end{enumerate}

\begin{enumerate}
\def\labelenumi{\alph{enumi})}
\item
  Show that the sets \(V_{u}\) are neighbourhoods of o in A = C(X; R) with respect to a Hausdorff topology C compatible with the ring structure of A.
\item
  The topology induced by \(\mathcal{C}\) on \(\mathcal{C}^{\infty}(\mathbf{X};\mathbf{R})\) is finer than the topology of uniform convergence in \(\mathbf{X}\) . These two topologies coincide if and only if \(\mathbf{X}\) is pseudo-compact (Chapter IX, § 1, Exercise 21). If \(\mathbf{X}\) is not pseudo-compact, the topology induced by \(\mathcal{C}\) on the set of constant functions is the discrete topology, so that \(\mathcal{C}\) is not compatible with the \(\mathbf{R}\) -vector space structure of \(\mathbf{A}\) .
\item
  Show that A is a Gelfand ring with respect to the topology \(\mathcal{C}\) (Chapter III, § 6, Exercise 11).
\end{enumerate}

¶ 15) Let X be a completely regular space and let βX be its Stone-Cech compactification. Endow the ring A = C(X; R) with the topology C defined in Exercise 14. For each f ∈ A, let V(f) be the closure in βX of the subset \(\overline{f}^{-1}(\mathrm{o})\) of X. For each subset M of A, let V(M) denote ∩f ∈ M V(f). For each subset Y of βX, let ℑ(Y) denote the set of all f ∈ A such that Y ⊂ V(f), and write ℑ(x) in place of ℑ(\{x\}). We have V(f) = ∅ if and only if f is a unit in A.

\begin{enumerate}
\def\labelenumi{\alph{enumi})}
\item
  If \(f, g\) are two elements of \(\mathbf{A}\) such that \(\overline{f}^{-1}(\mathrm{o}) \cap \overline{g}^{-1}(\mathrm{o}) = \emptyset\) , show that \(\mathrm{V}(f) \cap \mathrm{V}(g) = \emptyset\) . {[}Consider the function \(|f| / (|f| + |g|)\) . Deduce that, for each subset \(\mathbf{M}\) of \(\mathbf{A}\) , we have \(\mathrm{V}(\mathbf{M}) = \mathrm{V}(\mathfrak{a})\) , where \(\mathfrak{a}\) is the ideal generated by \(\mathbf{M}\) in \(\mathbf{A}'\) .
\item
  Show that, for each subset Y of \(\beta X\) , \(\Im(Y)\) is an ideal of A {[}use a){]} and that \(\mathrm{V}(\Im(\mathrm{Y}))\) is the closure \(\overline{Y}\) of Y in \(\beta X\) .
\item
  Let \(u \in \mathbf{A}\) be such that \(u(x) > 0\) for all \(x \in \mathbf{X}\) , and let \(g \in \mathbf{A}\) . Show that there exists \(f \in \mathbf{A}\) such that \(|f - g| \leqslant u\) and such that \(\mathbf{V}(f)\) is a neighbourhood of \(\mathbf{V}(g)\) . {[}Define \(f\) to be equal to \(g + u\) if \(g(x) + u(x) < 0\) , to \(g - u\) if \(g(x) - u(x) > 0\) , and to \(o\) otherwise; consider the function \(f' = \inf ((u + g)^+, (u - g)^-)\) and use \(a)\) .{]}
\item
  Let \(\mathfrak{a}\) be an ideal of \(\mathbf{A}\) . Show that if \(f\in\mathbf{A}\) is such that \(\mathbf{V}(f)\) is a neighbourhood of \(\mathbf{V}(\mathfrak{a})\) , then \(f\in\mathfrak{a}\) . {[}Note first that there exists \(g \in \mathfrak{a}\) such that \(\mathbf{V}(g)\) is contained in the interior of \(\mathbf{V}(f)\) , and then consider the function \(h\) defined on \(\mathbf{X}\) which is equal to \(f(x) / g(x)\) if \(f(x) \neq 0\) , and equal to \(0\) if \(f(x) = 0\) .{]}
\item
  Deduce from c) and d) that \(\Im(\mathrm{V}(\mathfrak{a})) = \overline{\mathfrak{a}}\) for every ideal \(\mathfrak{a}\) of A. {[}Show first that \(\mathrm{V}(\mathbf{M}) = \mathrm{V}(\overline{\mathbf{M}})\) for all subsets M of A, arguing by contradiction.{]} Deduce that, for each subset Y of \(\beta\mathrm{X}\) , \(\Im(\mathrm{Y})\) is a closed ideal of A. {[}Note that \(\Im(\mathrm{Y}) = \Im(\mathrm{V}(\Im(\mathrm{Y})))\) .{]}
\item
  Use e) to show that \(x \to \Im(x)\) is a bijection of \(\beta X\) onto the set of maximal ideals (necessarily closed) of A. An ideal of A is closed if and only if it is an intersection of maximal ideals. The ring A has no radical.
\end{enumerate}

¶ 16) We retain the hypotheses and notation of Exercise 15.

\begin{enumerate}
\def\labelenumi{\alph{enumi})}
\item
  Let a be a closed ideal in the ring A. Show that in the quotient ring A/a the set P of canonical images of functions \(f \geqslant 0\) of A is the set of elements \(\geqslant 0\) with respect to an ordering which makes A/a a lattice-ordered ring {[}observe, using Exercise 15 e), that if f and g belong to a, then so do \(|f| + |g|\) and every function \(h \in A\) such that \(|h| \leqslant |f|\) . The canonical image in A/a of the set of constant functions is isomorphic to R as an ordered field.
\item
  In particular, for each \(x \in \beta X\) , the field \(\mathrm{A}/\Im(x)\) is canonically endowed with an ordered field structure. The maximal ideal \(\Im(x)\) is said to be real if \(\mathrm{A}/\Im(x)\) is isomorphic to R, hyper-real otherwise. For \(\Im(x)\) to be hyper-real it is necessary and sufficient that \(x \in \beta X - X\) and that there is a unit \(f \in A\) such that \(\lim_{y \to x, y \in X} f(y) = 0\) . Deduce that all the maximal ideals of A are real if and only if X is pseudo-compact (Chapter IX, §1, Exercise 21).
\item
  For each \(x \in \beta\mathbf{X}\) , let \(\varphi_x\) be the canonical homomorphism \(\mathbf{A} \to \mathbf{A}/\Im(x)\) . For \(\varphi_x(f)\) to be \(\geqslant 0\) in \(\mathbf{A}/\Im(x)\) , it is necessary and sufficient that there should exist \(g \in \Im(x)\) such that \(f(y) \geqslant 0\) in \(\overline{g}^{-1}(0)\) {[}observe that the relation \(\varphi_x(f) \geqslant 0\) is equivalent to \(f \equiv |f|\pmod{\Im(x)}\) {]}. For \(\varphi_x(f)\) to be \(>0\) , it is necessary and sufficient that there should exist \(g \in \Im(x)\) such that \(f(y) > 0\) in \(\overline{g}^{-1}(0)\) . {[}Note that, if \(f \notin \Im(x)\) , there exists \(g \in \Im(x)\) such that \(\overline{f}^{-1}(0) \cap \overline{g}^{-1}(0) = \emptyset\) , by using Exercise 15.{]}
\item
  Show that if \(\Im(x)\) is hyper-real, the transcendence degree of the field \(\mathbf{A}/\Im(x)\) over \(\mathbf{R}\) is at least equal to \(\mathfrak{c}=\operatorname{Card}(\mathbf{R})\) . {[}Let \(f>0\) be a unit of \(\mathbf{A}\) such that \(\varphi_x(f)=u\) is infinitely large with respect to \(\mathbf{R}\) . Show that as \(r\) runs through the elements of a base of \(\mathbf{R}\) over \(\mathbf{Q}\) (consisting of numbers \(>0\) ), the elements \(\varphi_x(f^r)\) of \(\mathbf{A}/\Im(x)\) are algebraically independent.{]}
\end{enumerate}

* e) For each point \(a = (a_{1}, \ldots, a_{n}) \in \mathbb{R}^{n}\) , let \(f_{a}(\mathbf{X})\) denote the polynomial \(\mathbf{X}^{n} + a_{1}\mathbf{X}^{n-1} + \cdots + a_{n}\) . For each real number \(\xi\) , let \(\nu(\xi)\) be the sum of the multiplicities of the zeros \(z\) of \(f_{a}\) in \(\mathbf{C}\) such that \(\mathrm{R}(z) = \xi\) ; and for each integer \(k\) such that \(\mathrm{i} \leqslant k \leqslant n\) , let \(\rho_{k}(a)\) be the smallest real number \(\xi\) such that \(\sum_{\eta \leqslant \xi} \nu(\eta) \geqslant k\) . Show that each of the functions \(\rho_{k}\) is continuous on \(\mathbb{R}^{n}\) (use Rouché's theorem).

\begin{enumerate}
\def\labelenumi{\alph{enumi})}
\setcounter{enumi}{5}
\item
  Show that, for each \(x \in \beta X\) , the ordered field \(A/\Im(x)\) is real-closed. {[}Let \(F(X) = X^{n} + f_{1}X^{n-1} + \cdots + f_{n}\) be a polynomial of odd degree with coefficients in A; the functions \(g_{k}(y) = \rho_{k}(f_{1}(y), \ldots, f_{n}(y))\) , where the \(\rho_{k}\) are the functions defined in e), are continuous on X, and for each \(y \in X\) there is an index k such that \(F(g_{k}(y)) = 0\) ; deduce that the product \(F(g_{1}) \cdots F(g_{n})\) belongs to \(\Im(x).]\) *
\item
  Let X be an infinite discrete space and let \(y \to F_y\) be a bijection of X onto the set of finite subsets of X; for each \(z \in X\) let \(M_z\) be the set of all \(y \in X\) such that \(z \in F_y\) ; these sets form a base of a filter F on X. Let x be a point of \(\beta X\) such that the corresponding ultrafilter on X (Chapter I, §9, Exercise 27) is finer than F. Let B be a subset of A such that Card (B) ≤ Card (X), and let \(y \to g_y\) be a surjective mapping \(X \to B\) . For each \(z \in X\) put \(f(z) = \mathrm{i} + \sup_{y \in \mathbb{F}_z} g_y(z)\) . Show that \(f(z) > g_y(z)\) for all \(z \in M_y\) , and deduce that \(\varphi_x(f) > \varphi_x(g_y)\) for all \(y \in Y\) {[}use c){]}. Conclude that Card \(A/\Im(x)\) \textgreater{} Card X.
\end{enumerate}

¶ 17) We retain the hypotheses and notation of Exercise 15.

\begin{enumerate}
\def\labelenumi{\alph{enumi})}
\item
  A maximal ideal \(\Im(x)\) is real if and only if, for every infinite sequence \((g_n)\) of elements of \(\Im(x)\) , the intersection of the sets \(\overline{g_n}^{-1}(0)\) is not empty. {[}Use Exercise 16 b); to show that the condition is necessary show that, if \(\bigcap_{n}\overline{g_{n}}^{-1}(0) = \emptyset\) , the function \(f(y) = \sum_{n = 0}^{\infty}\inf (g_n(y),2^{-n})\) is continuous, is a unit in A and tends to o as \(y\) tends to \(x\) while remaining in X; note that \(f(y)\leqslant 2^{-m}\) for \(y\in \bigcap_{k\leqslant m}\overline{g_k}^{-1}(0)\) , and use Exercise 15 a) and the fact that the V(g), where \(g\in \Im(x)\) , form a fundamental system of neighbourhoods of \(x\) in \(\beta \mathbf{X}.\) {]}
\item
  X is said to be real-compact if the only real maximal ideals \(\Im(x)\) are those for which \(x \in \mathbf{X}\) . Show that every completely regular Lindelöf space (Chapter I, § 9, Exercise 15) is real-compact {[}use \(a\) {]}. A completely regular pseudo-compact space is not real-compact unless it is compact.
\item
  Let \(\nu\mathbf{X}\) be the set of all \(x\in \beta \mathbf{X}\) such that \(\Im (x)\) is real. Show that every continuous real-valued function \(f\in \mathbf{A}\) can be extended by continuity to \(\nu\mathbf{X}\) , so that \(\mathcal{C}(\mathbf{X};\mathbf{R})\) and \(\mathcal{C}(\upsilon\mathbf{X};\mathbf{R})\) can be canonically identified. Show that the subspace \(\nu\mathbf{X}\) of \(\beta\mathbf{X}\) is real-compact {[}note that if \(f\) is a unit in \(\mathcal{C}(\mathbf{X};\mathbf{R})\) , then its continuous extension to \(\nu\mathbf{X}\) is a unit in \(\mathcal{C}(\upsilon\mathbf{X};\mathbf{R})\) and that \(\beta(\upsilon\mathbf{X}) = \beta\mathbf{X}]\) . \(\nu\mathbf{X}\) is called the real-compactification of \(\mathbf{X}\) .
\item
  Show that \(\nu\mathbf{X}\) is the completion of \(\mathbf{X}\) with respect to the coarsest uniformity on \(\mathbf{X}\) for which all the functions \(f\in \mathbf{A}\) are uniformly continuous. (Note that a Cauchy filter \(\mathfrak{F}\) on \(\mathbf{X}\) with respect to this uniformity converges to a point \(x\in \beta \mathbf{X}\) ; if \(x\notin \mathfrak{vX}\) , there would be a function \(f\in \mathbf{A}\) which was unbounded in a set of F.{]}
\item
  Deduce from d) that if \(u: X \to Y\) is any continuous mapping of a completely regular space X into a real-compact space Y, then u can be extended by continuity to a mapping \(\upsilon X \to Y\) . The space \(\nu X\) is therefore the solution of the universal mapping problem in which \(\Sigma\) is the structure of a real-compact space, and the \(\alpha\) -mappings and morphisms are continuous mappings (Set Theory, Chapter IV, § 3, no. 1).
\end{enumerate}

\(f\) ) Show that every closed subspace of a real-compact space is real-compact, that every product of real-compact spaces is real-compact, and that in a Hausdorff space any intersection of real-compact subspaces is real-compact. {[}For example, to show that a closed subspace \(\mathbf{X}\) of a real-compact space \(\mathbf{Y}\) is real-compact, consider the extension \(\vec{j}\) to \(\nu\mathbf{X}\) of the canonical injection \(j:\mathbf{X}\to\mathbf{Y}\) , and note that \(j^{-1}(\mathbf{X})=\mathbf{X}\) ; deduce that \(\mathbf{X}\) is closed in \(\nu\mathbf{X}\) . The proofs are similar in the other two cases.{]}

\begin{enumerate}
\def\labelenumi{\alph{enumi})}
\setcounter{enumi}{6}
\tightlist
\item
  Deduce from b), d) and f) that a completely regular space is real-compact if and only if it is homeomorphic to a closed subspace of a product \(R^{I}\) .
\end{enumerate}

\begin{enumerate}
\def\labelenumi{\arabic{enumi})}
\setcounter{enumi}{17}
\tightlist
\item
  Let X be a real-compact space.
\end{enumerate}

\begin{enumerate}
\def\labelenumi{\alph{enumi})}
\item
  Show that, for each non-zero homomorphism \(u: \mathcal{C}(\mathbf{X}; \mathbf{R}) \to \mathbf{R}\) , there exists a unique point \(x \in \mathbf{X}\) such that \(u(f) = f(x)\) for all \(f\) in \(\mathcal{C}(\mathbf{X}; \mathbf{R})\) .
\item
  Let Y be another real-compact space. Show that, for every R-algebra homomorphism \(v: \mathcal{C}(Y; \mathbf{R}) \to \mathcal{C}(X; \mathbf{R})\) which maps identity element to identity element, there exists a unique continuous mapping \(w: X \to Y\) such that \(v(g) = g \circ w\) for all \(g \in \mathcal{C}(Y; \mathbf{R})\) . In particular, \(\mathcal{C}(X; \mathbf{R})\) and \(\mathcal{C}(Y; \mathbf{R})\) are isomorphic if and only if X and Y are homeomorphic.
\end{enumerate}

¶ 19) Let X be a compact space, and let A be the normed C-algebra C(X; C). For each closed ideal a of A, show that the relation \(f \in a\) implies \(\overline{f} \in a\) (note that \(\overline{f}\) is the uniform limit of functions of the form \(g f\) ). Extend the results of Exercise 13 to the algebra A.

¶ 20) a) Let X be a compact space, let A be an R-algebra of finite rank with an identity element, and endow A with the topology defined in Chapter VI, § 1, no. 5. Let B be a subalgebra of the R-algebra C (X; A), and suppose that, for each ε \textgreater{} o, each pair of distinct points x, y of X and each pair of elements u, v of A, there exists f ∈ B such that \textbar\textbar f(x) - u\textbar\textbar{} ≤ ε and \textbar\textbar f(y) - v\textbar\textbar{} ≤ ε (\textbar\textbar{} \textbar\textbar{} being any norm defining the topology of A); and suppose, further, that if u and v belong to R, there is a function f ∈ B satisfying these conditions and taking its values in R. Show that, under these hypotheses, B is dense in C (X; A) with respect to the topology of uniform convergence. {[}Show first that every function of C (X; R) can be uniformly approximated by functions of B and then that the same is true of every constant a ∈ A, by using the hypotheses and a partition of unity.{]}

\begin{enumerate}
\def\labelenumi{\alph{enumi})}
\setcounter{enumi}{1}
\tightlist
\item
  Take A = H, the division ring of quaternions over R. Let B be a subalgebra of the R-algebra \(\mathcal{C}(\mathbf{X}; \mathbf{H})\) such that: (i) for each \(f \in B\) , the conjugate \(\overline{f}\) of f {[}defined by \(\overline{f}(x) = \overline{f(x)}\) for all \(x \in X\) {]} belongs to B; (ii) for each \(\varepsilon > 0\) , each pair of distinct points x, y of X and each pair of elements u, v of H, there exists \(f \in B\) such that \(||f(x) - u|| \leqslant \varepsilon\) and \(||f(y) - v|| \leqslant \varepsilon\) . Show that B is dense in \(\mathcal{C}(\mathbf{X}; \mathbf{H})\) {[}use a){]}. c) Extend the results of Exercise 13 to the algebra \(\mathcal{C}(\mathbf{X}; \mathbf{H})\) , and show in particular that every closed ideal of this algebra is two-sided {[}argue as in Exercise 19, using b){]}.
\end{enumerate}

¶ 21) a) Let X be a totally disconnected compact space, and let A be a topological ring. Endow the ring C (X; A) with the topology of uniform convergence, which is compatible with the ring structure. Let a be a left ideal of C (X; A) and, for each \(x \in X\) , let \(\Im(x)\) be the closure in A of the set of \(f(x)\) for \(f \in a\) ; \(\Im(x)\) is a closed left ideal of A. Show that \(\overline{a}\) is identical with the set of all \(f \in C(X; A)\) such that \(f(x) \in \Im(x)\) for all \(x \in X\) (note that, for each \(x \in X\) , the neighbourhoods of x which are both open and closed form a fundamental system of neighbourhoods of x).

\begin{enumerate}
\def\labelenumi{\alph{enumi})}
\setcounter{enumi}{1}
\tightlist
\item
  Suppose in addition that A has a fundamental system of neighbourhoods of o which are two-sided ideals. Let B be a subring of \(\mathcal{C}(\mathbf{X};\mathbf{A})\) which contains the constants and separates the points of X. Show that B is dense in \(\mathcal{C}(\mathbf{X};\mathbf{A})\) . {[}For each closed subset F of X, each \(x\notin F\) and each neighbourhood U of o in A, show that there exists \(f\in B\) such that \(f(x)=1\) and \(f(y)\in U\) for all \(y\in F\) .{]}
\end{enumerate}

